\appendix

\section{Sign discussion}
\label{AppendixSign}

In this appendix, we fix the sign indeterminacy present in the proof of Theorem~\ref{ThmRepeatedMain} above.
Precisely, we prove the following.

\begin{Proposition}
\label{SignProposition}
For $d\geq 5$, the transgression
\begin{equation*}
\tau(p_1) \in H^3(\L B\SO(d), \Z) = H^3(B\L\SO(d), \Z)
\end{equation*}
of the universal first Pontrjagin class equals two times the pullback of the universal Dixmier--Douady class $\DD \in H^3(B\Aut(A_d), \Z)$ along the map on classifying spaces induced by the composition
\begin{equation*}
\begin{tikzcd}
\tilde{\theta} \colon\ \L\SO(d) \ar[r, "j"] & \O_{\mathrm{res}}\big(H^d\big) \ar[r, "\theta"] & \Aut(A_d).
\end{tikzcd}
\end{equation*}
\end{Proposition}

For definiteness, we emphasize that we use the (standard) convention for Pontrjagin classes that for any complex line bundle $L$ with underlying real bundle $L_\R$, we have $p_1(L_\R) = c_1(L)^2$.

We recall the definition of the universal Dixmier--Douady class $\DD$.
To begin with, recall that for a topological group $G$ with classifying space fibration $G \to EG \to BG$, there is a~homomorphism
\begin{equation*}
%\label{TildeTransgressionDef}
 \tilde{\tau} \colon\ H^k(BG, R) \longrightarrow H^{k-1}(G, R),
\end{equation*}
natural in $G$, which is called transgression and should not be confused with the notion of transgression discussed in Section~\ref{SectionTransgression}.\footnote{Both this notion of transgression and the one discussed in Section~\ref{SectionTransgression} are special cases of the more general notion of transgression in a fiber bundle \cite{Borel}.}
In the case of $G = \Aut(A_d)$ and $k=3$, the transgression homomorphism
\begin{equation}
\label{TildeTransgression}
\tilde{\tau} \colon\ H^3(B\Aut(A_d), \Z) \longrightarrow H^2(\Aut(A_d), \Z)	
\end{equation}
is an isomorphism, and the universal Dixmier--Douady class $\DD$ is by definition the class that is sent to the first Chern class $c_1(G_d) \in H^2(\Aut(A_d), \Z)$ of the canonical $\U(1)$-bundle $G_d$ over~$\Aut(A_d)$.




As the classes we are interested in are not torsion, we may work over real coefficients.
We consider the diagram
\begin{equation*}
\begin{tikzcd}[column sep=0.8cm]
H^3(B\Aut(A_d), \R)
	\ar[r, "B\theta^*", "\cong"']
	\ar[d, "\tilde{\tau}"]
&
H^3\big(B\O_{\res}\big(H^d\big), \R\big)
	\ar[r, "Bj^*", "\cong"']
	\ar[d, "\tilde{\tau}"]
&
H^3(B\L\SO(d), \R)
	\ar[d, "\tilde{\tau}"]
&
H^4(B\SO(d), \R)
\ar[l, "\tau"', "\cong"]
	\ar[d, "\tilde{\tau}"]
\\
H^2(\Aut(A_d), \R)
	\ar[r, "\theta^*"]
&
H^2\big(\O_{\res}\big(H^d\big), \R\big)
	\ar[r, "j^*", "\cong"']
&
H^2(\L \SO(d), \R)
&
H^3(\SO(d), \R),
\ar[l, "\tau"']
\end{tikzcd}
\end{equation*}
which commutes by naturality of the transgression maps \eqref{TildeTransgression}.
By the results of Section~\ref{sec4} and the universal coefficient theorem, all groups in this diagram are isomorphic to $\R$ and all maps are isomorphisms.

In a first step, we observe that the pullback of the bundle $G_d$ under the map $\theta$ is precisely the implementer bundle $\Imp \to \O_{\mathrm{res}}\big(H^d\big)$, so that the statement of Proposition~\ref{SignProposition} is equivalent to the equality
\begin{equation*}
%\label{LeftToShow}
	2\cdot j^*c_1(\Imp) = \tau (\tilde{\tau}(p_1)).
\end{equation*}

Since the three groups on the right are all Fr\'echet Lie groups, we may work with de Rham cohomology instead of singular cohomology.
Here the transgression homomorphisms $\tilde{\tau}$ may be described as follows.
Let $\alpha \in H^k_{\mathrm{dR}}(BG)$ and let $\pi^*\alpha \in H^k_{\mathrm{dR}}(EG)$ be its pullback.
Then since~$EG$ is contractible, $\pi^* \alpha = d \beta$ for some $\beta \in \Omega^k(EG)$.
The transgression of $\alpha$ is then defined by
\begin{equation}
\label{DefTransgressionTilde}
	\tilde{\tau}(\alpha) = \iota^*\beta,
\end{equation}
where $\iota \colon \SO(d) \to E\SO(d)$ is the inclusion of a fiber.

The de Rham cohomology groups $H^k_{\mathrm{dR}}(\SO(d))$ are best understood using the isomorphism with the Lie algebra cohomology group $H^k(\mathfrak{so}(d), \R)$, which identifies a Lie algebra $k$-cocycle $\alpha$ with the left-invariant 3-form $\overline{\alpha}$ on $\SO(d)$ that coincides with $\alpha$ at the identity.
We now have the following general statement.

\begin{Lemma}
\label{SecondLemmaA}
We have $\tilde{\tau}(p_1) = \overline{\sigma}$, the left invariant $3$-form corresponding to the Lie algebra cocycle
\begin{equation*}
 \sigma(x, y, z) = \frac{1}{8 \pi^2} \langle x, [y, z] \rangle, \qquad x, y, z \in \mathfrak{so}(d).	
\end{equation*}
\end{Lemma}

The proof uses the theory of Chern and Simons \cite{ChernSimons}, which we recall now.
Let $\omega$ be a connection 1-form on a principal $G$-bundle $E$ over $B$ with curvature $\Omega$.
Let $P \in \mathrm{Sym}^2(\mathfrak{g}^*)^G$ be an invariant polynomial on $\mathfrak{g}$.
Then the corresponding Chern--Simons form is $\mathrm{T}P(\omega) \in H^3(E, \R)$ defined by
\begin{equation*}
 \mathrm{T}P(A) = P(\omega \wedge \Omega) - \frac{1}{6} P(\omega \wedge [\omega, \omega]),
\end{equation*}
see \cite[formula~(3.5)]{ChernSimons}.
Denoting by $\iota \colon G \to E$ the inclusion of a fiber (for some fixed base point~$e \in E$), one uses that the curvature form $\Omega$ is horizontal, so that $\iota^*\Omega = 0$, while $\iota^*\omega = \omega_G$, the Maurer--Cartan form of $G$.
Hence
\begin{equation*}
%\label{PullbackTP}
\iota^*\mathrm{T}P(\omega) = -\frac{1}{6} P(\omega_G \wedge [\omega_G, \omega_G]),
\end{equation*}
which differs from the formula (3.11) in \cite{ChernSimons} by a factor of 2.
Going through the conventions used in \cite{ChernSimons} for the wedge product and commutator of $\mathfrak{g}$-valued differential forms (see \cite[p.~50]{ChernSimons}), one obtains that $\iota^* \mathrm{T}P(\omega)$ is the left-invariant form corresponding to the Lie algebra cocycle
\begin{equation}
\label{DefinitionAlpha}
 \alpha(x, y, z) =	- \frac{1}{3} P(x \otimes [y, z] + y \otimes [z, x] + z \otimes [x, y]).
\end{equation}

\begin{proof}
We take $G = \SO(d)$ and $E = E\SO(d)$, $B = B\SO(d)$, the universal bundles.
Choosing models for these that are infinite-dimensional Fr\'echet manifolds (for example, the infinite Grassmannian and Stiefel manifold), we may choose a connection 1-form $\omega \in H^1(E\SO(d), \mathfrak{so}(d))$.
Setting $P(X \otimes Y) = -\tr(XY)/8\pi^2$, \cite[Proposition~3.2]{ChernSimons} states that
\begin{equation*}
 d \mathrm{T}P(\omega) = P(\Omega \otimes \Omega) = -\frac{1}{8\pi^2}\tr\big(\Omega^2\big).
\end{equation*}
By the usual Chern--Weil formulas for Pontrjagin classes, we see that this equals the pullback~$\pi^*p_1$ of the de Rham representative of the universal first Pontrjagin class along the bundle projection~$\pi\colon E\SO(d) \to B\SO(d)$.

Using the description \eqref{DefTransgressionTilde} of the transgression homomorphism, we see that $\tilde{\tau}(p_1) = \iota^*\mathrm{T}P(\omega)$.
As discussed above, the pullback $\iota^*\mathrm{T}P(\omega)$ is the left invariant differential form that corresponds to the Lie algebra cocycle $\alpha$ given by \eqref{DefinitionAlpha}.
For our particular choice of $P$, we get
\[
\alpha(x, y, z) = \left(-\frac{1}{3}\right) \cdot\left(-\frac{1}{8\pi^2}\right) \tr(x [y, z] + y [z, x] + z [x, y])
= \sigma(x, y, z).\tag*{\qed}
\]\renewcommand{\qed}{}
\end{proof}



\begin{Lemma}
\label{LemmaA3}
We have $\tau(\overline{\sigma}) = -\overline{\omega}/2 \pi$, where $\overline{\omega}$ is the left-invariant $2$-form on $L\SO(d)$ corresponding to the Lie algebra cocycle
\begin{equation*}
 \omega(X, Y) = \frac{1}{2\pi}\int_{S^1} \tr(X(t)Y^\prime(t)), \qquad X, Y \in \L\mathfrak{so}(d).	
\end{equation*}
\end{Lemma}

This result can be found as \cite[Proposition~4.4.4]{PressleySegal}, but with incorrect prefactors and an incomplete proof, which is why we repeat the proof below.

\begin{proof}
Tangent vectors $X$, $Y$ at $\gamma \in L\SO(d)$ can be identified with those elements of $\L \mathrm{Mat}_{d\times d}$ such that $\tilde{X} \in L\mathfrak{so}(d)$, where $\tilde{X}(t) = \gamma(t)^{-1}X(t)$. We now calculate
 \begin{align*}
 \tau(\overline{\sigma})_\gamma(X, Y)
 &= \int_{S^1} \overline{\sigma}(\gamma^\prime(t), X(t), Y(t)) {\rm d}t= \int_{S^1} \sigma\big( \gamma(t)^{-1}\gamma^\prime(t), \gamma(t)^{-1} X(t), \gamma(t)^{-1}Y(t)\big) {\rm d}t\\
 &= \frac{1}{8 \pi^2}\int_{S^1}\tr\big(\gamma(t)^{-1}\gamma^\prime(t) \big[\tilde{X}(t), \tilde{Y}(t)\big]\big) {\rm d}t.
 \end{align*}
 Let $\beta \in \Omega^1(\L\SO(d))$ be given by
 \begin{equation*}
 \beta_\gamma(X) = \int_{S^1} \tr \big( 	\gamma(t)^{-1}\gamma^\prime(t) \tilde{X}(t)\big) {\rm d}t.
 \end{equation*}
For the exterior derivative of $\beta$, we find
 \begin{align*}
 d\beta_\gamma(X, Y) &= \partial_X \{ \beta(Y)\}_\gamma - \partial_Y \{ \beta(X)\}_\gamma - \beta_\gamma([X, Y]) \\
 &= \int_{S^1} \big( \tr \big(\tilde{X}^\prime(t) \tilde{Y}(t)\big) - \tr \big(\tilde{Y}^\prime(t) \tilde{X}(t)\big) - \tr\big(\gamma(t)^{-1}\gamma^\prime(t) [\tilde{X}(t), \tilde{Y}(t)]\big)\big) {\rm d}t
 \end{align*}
 for suitable extensions of $X$ and $Y$ to vector fields on $\L\SO(d)$.
 Integrating by parts, we see that ${\rm d} \beta = - 2 \cdot 2\pi \cdot \overline{\omega} - 8 \pi^2 \cdot \tau(\overline{\sigma})$, hence $2 \pi \cdot \tau(\overline{\sigma})$ and $-\overline{\omega}$ are cohomologous.
\end{proof}

Let $\tilde{G} \to G$ be a central $\U(1)$-extension of a Fr\'echet Lie groups, inducing a Lie algebra homomorphism $\tilde{\mathfrak{g}} \to \mathfrak{g}$.
Choosing a linear section of the Lie algebra homomorphism $\tilde{\mathfrak{g}} \to \mathfrak{g}$ gives an identification $\tilde{\mathfrak{g}} = \mathfrak{g} \oplus \R$.
With respect to this choice, the Lie bracket of $\tilde{\mathfrak{g}}$ is given by
\begin{equation*}
 	[(X, \lambda), (Y, \mu)] = ([X, Y], \Omega(X, Y))
\end{equation*}
for some continuous Lie algebra cocycle 2-cocycle $\Omega$, which represents a class in $H^2_c(\mathfrak{g}, \R)$.
The first Chern class of the principal $\U(1)$-bundle $\tilde{G} \to G$ is then given by
\begin{equation}
\label{ChernClassVsCocycle}
 c_1(\tilde{G}) = \frac{1}{2\pi} \overline{\Omega}, 	
\end{equation}
where $\overline{\Omega}$ denotes the associated left-invariant 2-form, see for example \cite[Proposition~4.5.6]{PressleySegal}.

In the case that $G = \O_{\res}(H, [L])$ and $\tilde{G} = \Imp_L$ for some real Hilbert space with a Lagrangian~$L \subset H_\C$, the relevant Lie algebra cocycle has been computed in several places; see, e.g., \cite[Theorem~6.10]{Araki1}, \cite[Theorem~10.2]{NeebSemibounded} or \cite[Theorem~6]{Ottesen}.
The result is
\begin{equation}
\label{GroupCocycleImp}
 \Omega(X, Y) = \frac{1}{8} \tr(J [J, X] [J, Y]), \qquad X, Y \in \mathfrak{o}_{\res}\big(H^d\big),
\end{equation}
where $\Gamma = {\rm i}(P_{L} - P_{\overline{L}})$ is the complex structure determined by the Lagrangian $L$.

When trying to apply these results to the specific Hilbert space $H^d$ defined in \eqref{DefinitionHd}, we face the difficulty that the subspace $L^d \subset H^d$ defined in \eqref{SubLagrangianLd} is only a sub-Lagrangian.
We deal with this issue as follows:
\begin{enumerate}\itemsep=0pt
\item[(i)]
If $d$ is even we let $K = \big(L^d \oplus \overline{L}^d\big)^\perp$ be the even-dimensional subspace of constant functions and choose a Lagrangian $L_0 \subset K$.
Then $L = L^d + L_0 \subset H^d_\C$ is a Lagrangian.
\item[(ii)]
If $d$ is odd, we let $K = \big(L^d \oplus \overline{L}^d\big)^\perp \oplus \C$ and again choose a Lagrangian $L_0 \subset K$.
Then~${L = L^d + L_0}$ is a Lagrangian in $H^d_\C \oplus \C$.
By definition, the implementer bundle over $\O_{\res}\big(H^d\big)$ is the restriction of the implementer bundle over $\O_{\res}\big(H^d \oplus \R\big)$, hence the group cocycle associated to this extension is the restriction of the cocycle \eqref{GroupCocycleImp} to~$\mathrm{o}_{\res}\big(H^d\big) \subset \mathfrak{o}_{\res}\big(H^d \oplus \R\big)$.
Given an element $g \in L\SO(d)$, the element $j(g)$ acts on~$H^d \oplus \R$ through multiplication by $g$ on the first summand and the identity on the second, and consequently, for $X \in \L \mathfrak{so}(d)$, the operator $j_*X$ acts by multiplication with $X$ on the first summand and by zero on the second.
\end{enumerate}
To have a uniform notation, we write $H$ for either the Hilbert space $H^d$ or for $H^d \oplus \R$ in the case that $d$ is odd and let $L \subset H$ be the Lagrangian described above.

\begin{Lemma}\label{LemmaA4}
 We have $2 \cdot j^*c_1(\Imp) = - \overline{\omega}/2\pi$.
\end{Lemma}

A similar computation, for the Lagrangian \eqref{OtherLagrangian} instead of $L$, can be found in \cite{KristelWaldorf1}; unfortunately, the result is off by a factor of $\pm {\rm i}$.
Proposition~6.7.1 in \cite{PressleySegal} is an analogous result for the basic central extension of restricted unitary group $\U_{\res}(H)$ of a polarized complex Hilbert space~$H$.

\begin{proof}We will establish the cocycle identity
\begin{equation}\label{EquationOfCocycles}
 2 \cdot j^*\Omega = -\omega,	
\end{equation}
which gives the result by \eqref{ChernClassVsCocycle}.
By continuity and bilinearity, it suffices to verify that both Lie algebra cocycles evaluate identically on the specific Lie algebra elements $X,Y \in L\mathfrak{so}(d)$ of the form
\begin{equation*}
	X(t) = a {\rm e}^{-{\rm i}kt}, \qquad Y(t) = b {\rm e}^{-{\rm i}\ell t}, \qquad a, b \in \mathfrak{so}(d), \qquad k, \ell \in \Z.
\end{equation*}
On these elements, the right-hand side of \eqref{EquationOfCocycles} is given by
\begin{align}
 \omega(X, Y) &= \frac{1}{2\pi}\int_0^{2\pi} \tr(X(t) Y^\prime(t)) {\rm d}t
 \nonumber\\
 &= -\frac{{\rm i}\ell}{2\pi}\int_0^{2\pi} \tr(ab) {\rm e}^{-{\rm i}(k+\ell)t} {\rm d}t
 = \begin{cases} -{\rm i}\ell \cdot \tr(ab) & k+\ell = 0, \\ 0 & \text{otherwise}.
 \end{cases}\label{ToCompareWithOmega}
\end{align}

To calculate the left-hand side of \eqref{EquationOfCocycles}, we write
	\begin{equation*}
 j_*X =
 \begin{pmatrix}
 	x^\prime & \overline{x} \\ x & \overline{x}^\prime
 \end{pmatrix},	
 \qquad
 j_*Y =
 \begin{pmatrix}
 	y^\prime & \overline{y} \\ y & \overline{y}^\prime
 \end{pmatrix}	
\end{equation*}
with respect to the decomposition $H_\C = L \oplus \overline{L}$.
In other words, $x^\prime = P_LXP_L$ and $x = P_{\overline{L}}XP_L$ and similarly for $y^\prime$ and $y$, while $\overline{x}$ denotes the conjugation of $x$ by the real structure of $H_\C$.
Then since $j_*X$ and $j_*Y$ are restricted, the off-diagonal entries $x$ and $y$ are Hilbert--Schmidt operators and a straightforward calculation gives
\begin{equation}
\label{ReformulatedOmega}
 \Omega(j_*X, j_*Y) = \frac{{\rm i}}{2} \tr(\overline{x} y - \overline{y}x).	
\end{equation}

Write $V_n = \big\{\xi \otimes {\rm e}^{\rm int} \mid \xi \in \C^d\big\} \subset H^d_\C$.
We observe that $j_*X$ sends $V_n$ to $V_{n+k}$ and $j_*Xj_*Y$ sends $V_n$ to $V_{n-k-\ell}$.
On the other hand, $P_L$ and $P_{\overline{L}}$ preserve the subspaces $V_n$ for $n \neq 0$ and send~$V_0$ to $K$ (where $K = V_0$ if $d$ is even, while $K = V_0 \oplus \C$ if $d$ is odd).
We conclude that both~$\overline{x} y$ and $x \overline{y}$ send $V_n$ to $V_{n-k-\ell}$, respectively $V_{n-k-\ell} \oplus \C$ if $d$ is odd.
Choosing an orthonormal basis adapted to the decomposition $H_\C = \bigoplus_{n \in \Z} V_n$, respectively $H_\C = (\bigoplus_{n \in \Z} V_n) \oplus \C$ to calculate the trace on the right-hand side of \eqref{ReformulatedOmega}, we see that the result can be non-zero only if $k+\ell = 0$.

So suppose now that $k = - \ell$ and let $\xi \otimes {\rm e}^{{\rm int}} \in V_n^\C$, $n \geq 1$. Then we have
\begin{align*}
 \overline{x}y(\xi\otimes \chi_n) &=
 \begin{cases}
 ab\xi \otimes \chi_{n} & 1 \leq n \leq \ell -1, \\
 a\overline{P}_0b\xi \otimes \chi_{n} & n = \ell,\\
 0 & \text{otherwise},
 \end{cases}
 \\
 \overline{y}x(\xi \otimes \chi_n) &=
 \begin{cases}
 ba\xi \otimes \chi_{n}, & 1 \leq n -1 \leq -\ell -1, \\
 b\overline{P}_0a\xi \otimes \chi_{n}, & n = - \ell,\\
 0, & \text{otherwise},
 \end{cases}	
\end{align*}
where $P_0$ is the orthogonal projection onto $L_0$ in $K$.
For $\xi \otimes 1 \in V_0$, we find
\begin{align*}
 \overline{x}y(\xi\otimes 1) &=
 \begin{cases}
 P_0ab\xi \otimes 1, & \ell > 0, \\
 P_0a\overline{P}_0b\xi \otimes 1, & \ell = 0,\\
 0, & \text{otherwise},
 \end{cases}
 \\
 \overline{y}x(\xi \otimes 1) &=
 \begin{cases}
 P_0ba\xi \otimes 1, & \ell < 0, \\
 P_0b\overline{P}_0a\xi \otimes 1, & \ell = 0,\\
 0, & \text{otherwise}.
 \end{cases}	
\end{align*}
Concluding, if $\ell >0$, we obtain
\begin{equation*}
 \Omega(j_*X, j_*Y) = \frac{{\rm i}}{2} \tr(\overline{x}y) = \frac{{\rm i}}{2} \left(\sum_{n=1}^{\ell-1}\tr(ab) + \tr\big(a\overline{P}_0 b\big) + \tr(P_0ab) \right) = \frac{{\rm i}\ell}{2}\tr(ab).
\end{equation*}
Here we used that
\[\tr\big(a\overline{P}_0 b\big) + \tr(P_0ab) = \tr(ab),\] which follows from the following calculation.
First, because $a$ and $b$ are real and skew-adjoint and $P_0$ is self-adjoint, we get
\begin{equation*}
 	\tr(aP_0b) = \tr(baP_0) = \overline{\tr((baP_0)^*)} = \overline{\tr(P_0 (-a)(-b))} = \tr\big(\overline{P}_0ab\big),
\end{equation*}
hence
\begin{equation*}
	\tr\big(a\overline{P}_0 b\big) + \tr(P_0ab) = \tr\big(\overline{P}_0 ab\big) + \tr(P_0ab) = \tr(ab),
\end{equation*}
using $\overline{P}_0 + P_0 = \id$.
If $\ell < 0$, we obtain in a similar fashion
\begin{equation*}
 \Omega(j_*X, j_*Y) = -\frac{{\rm i}}{2} \tr(\overline{y}x) = -\frac{{\rm i}}{2} \left(\sum_{n=1}^{|\ell|-1}\tr(ba) + \tr\big(b\overline{P}_0 a\big) + \tr(P_0ba) \right) = -\frac{{\rm i}|\ell|}{2}\tr(ab).
\end{equation*}
Finally, if $\ell=0$, then
\begin{equation*}
\begin{aligned}
 	\Omega(j_*X, j_*Y) = \frac{{\rm i}}{2} \tr(\overline{x}y - \overline{y}x)
 	&= \frac{{\rm i}}{2} \big(\tr\big(P_0 a\overline{P}_0b\big) - \tr\big(P_0 b\overline{P}_0a\big)\big)\\
 	&= \frac{{\rm i}}{2} \big(\tr(P_0 ab) - \tr(P_0 a{P}_0b) - \tr(P_0 ba) + \tr(P_0 bP_0a)\big) \\
 	&=\frac{{\rm i}}{2} \big(\tr(P_0 ab) - \tr(ba) + \tr\big(\overline{P}_0 ba\big) \big) \\
 	& = 0.
\end{aligned}
\end{equation*}
Comparing with \eqref{ToCompareWithOmega}, this establishes \eqref{EquationOfCocycles}.
\end{proof}
